%% kapa: exam=1 round=30 year=2019 period=1 form=A subject=civil_law subject_ko=민법 seq=18 number=18
%% source: Q-Net 감정평가사 기출문제 (cst00309) artlSeq=5205205 file=extracted/1008/5205205/2019년 제30회 감정평가사 1차시험 문제지/2019년 제30회 감정평가사 1차시험 문제지/2019년도 제30회 감정평가사 1차 1교시 A형.hwp
%% conversion: kordoc 4.24.51 -> tools/gen1.py (auto); review=proofread-sonnet; flags=none
\begin{question}{18}
기간에 관한 설명으로 옳지 않은 것은? (단, 기간 말일이 토요일 또는 공휴일인 경우는 고려하지 않음)
\begin{choices}
  \item 기간을 시, 분, 초로 정한 때에는 즉시로부터 기산한다.
  \item 채무의 이행기를 일, 주, 월 또는 연으로 정한 때에는 기간이 오전 0시로부터 시작하는 경우가 아닌 한, 기간의 초일을 산입하지 않는다.
  \item 기간을 일, 주, 월 또는 연으로 정한 때에는 기간말일의 종료로 기간이 만료한다.
  \item 연령을 계산하는 경우에는 출생일을 산입한다.
  \item 주, 월 또는 연의 처음부터 기간을 기산하지 아니한 때에는 최후의 주, 월 또는 연에서 그 기산일에 해당한 날로 기간이 만료한다.
\end{choices}
\end{question}
\answer{5}
